% !TeX root=../main.tex
\chapter{مروری بر مطالعات انجام شده}

% TODO ۲-۱ مقدمه فصل — یک پاراگراف کوتاه درباره‌ی چارچوب مرور: نظری (Newsvendor) + صنایع
% آنالوگ + کترینگ نهادی + جعبه‌ابزار فنی.
\section{مقدمه}

% TODO ۲-۲ تعاریف، اصول و مبانی نظری
% این بخش الزامی نیست ولی برای این پروژه ضروری است — دو خواسته‌ی صریح کارفرما اینجا
% جای می‌گیرند:
\section{تعاریف، اصول و مبانی نظری}

% TODO ۲-۲-۱ مسئله‌ی روزنامه‌فروش (Newsvendor) و زیان نامتقارن
% عدم‌تقارن هزینه‌ی کمبود Cu در برابر مازاد Co؛ کوانتایل بهینه، رابطه (۱):
% tau*_rho = Co/(Cu+Co) در فضای نرخ عدم‌دریافت؛ چرا فرمول اولیه‌ی سند مسئله وارونه بود
% (تصحیح مستند شده به‌عنوان دستاورد، نه اشتباه پنهان‌شده).
% منبع: doc/project-definition.md بند ۲-۲؛ ⚠️ ردیف ۳۶ doc/decision_log.md
\subsection{مسئله‌ی روزنامه‌فروش و زیان نامتقارن}

% TODO ۲-۲-۲ رگرسیون کوانتایل و تابع زیان Pinball
% رابطه (۲)؛ معادل‌سازی کمینه‌ساز Pinball با کوانتایل tau؛ تفاوت کوانتایل با میانگین
% (⚠️ تناقض ظاهری log_res — ردیف ۵۱)؛ تله‌ی تجمیع «جمع کوانتایل‌ها != کوانتایل جمع».
% منبع: doc/Literature_Review.md بند ۲-۲، ۲-۳-۱
\subsection{رگرسیون کوانتایل و پیش‌بینی احتمالاتی}

% TODO ۲-۲-۳ کالیبراسیون، پوشش و بیش‌پراکندگی
% پوشش حاشیه‌ای در برابر شرطی؛ Conformal/CQR و Mondrian؛ نقض فرض استقلال افراد در
% تجمیع سطح فرد (F08) و پیامدش برای کوانتایل خوش‌بینانه.
\subsection{کالیبراسیون و بیش‌پراکندگی}

% TODO ۲-۲-۴ اعتبارسنجی و نشت اطلاعاتی در داده‌ی زمانی
% چرا K-fold تصادفی معتبر نیست؛ rolling-origin؛ انواع نشت (هدف، زمانی، Target Encoding).
% منبع: doc/progress/06-*.md؛ doc/leakage_audit.md
\subsection{اعتبارسنجی و نشت اطلاعاتی در داده‌ی زمانی}

% TODO ۲-۲-۵ ابزارهای آماری و استنباطی به‌کاررفته
% جدول: هر آزمون، چه می‌سنجد، فرض‌ها. دوازده مورد: ADF/KPSS، Ljung-Box، ARCH-LM،
% Kruskal-Wallis، Mann-Whitney/Wilcoxon، Levene، Spearman/Kendall، Diebold-Mariano،
% بوت‌استرپ بلوکی دوبعدی و n_eff، Cliff's delta، تصحیح چندگانه BH-FDR.
% فقط تعریف اینجا؛ نتیجه‌ی هر آزمون در فصل ۳/۴ با ارجاع به همین بند.
\subsection{ابزارهای آماری و استنباطی به‌کاررفته}

% TODO ۲-۳ مروری بر ادبیات موضوع
% چهار محور: نظری Newsvendor · صنایع آنالوگ no-show (هواپیمایی/هتل/رستوران) ·
% کترینگ نهادی/دانشگاهی · جعبه‌ابزار فنی (رقابت M5 Uncertainty). جدول فشرده‌ی ۱۵ منبع.
% خلأ پژوهشی: نبود پیشینه‌ی مستقیم ایرانی.
% منبع: doc/Literature_Review.md (کل)
\section{مروری بر ادبیات موضوع}

% TODO ۲-۴ نتیجه‌گیری فصل
\section{نتیجه‌گیری}
